
%\chapter{Literature Review}

\section{Literature Review}
\label{literature}

The LRSP integrates three sets of decisions, facility location, 
vehicle routing and vehicle scheduling. In this chapter, we briefly 
review the literature for each component problem as well as the 
literature for related integrated problems. 

\subsection{Facility Location}

The first component of the LRS problem is where to locate the facilities 
and how to allocate the customers to these facilities. Determining where 
to locate facilities is an important problem that arises in the design of 
distribution networks. \citet{ReVelle} and \citet{Klose} provide recent 
surveys of the various types of location problems, formulations and solution 
algorithms. The facility location literature encompasses a wide variety of 
problems, which can be categorized in a variety of ways. \citet{Daskin} 
provides a taxonomy of facility location problems and describes methods 
categorizing location problems. Location problems can be grouped as planar 
problems, network problems or discrete problems based on the allowable 
location of the facilities and the customers. Problems in which customers 
and facilities can be located anywhere in the plane are planar location 
problems. Problems in which the customers and the candidate locations are 
restricted to a discrete set  are discrete location problems. Network 
location problems are a special case of discrete location problems in which 
the candidate facilities and the demand locations are restricted to the nodes 
of a network and travel is only possible via the arcs of network.
%[alternative to the previous]Problems in which the candidate facilities and the demand locations are 
%restricted to the nodes of a network and travel is only possible via the arcs of network are network 
%location problems. 

Most relevant to our study of the LRS problem are discrete location problems. 
\citet{Labbe} provide a survey of the literature related to discrete facility 
location problems including application areas, solution methods and model 
extensions. Early work on discrete location problems started with 
\citet{Hakimi64, Hakimi65} who introduced the $p$-median problem on a network, 
which is known to be NP-hard (\citet{Kariv}). The objective is to locate $p$ 
facilities on a network with demand nodes so as to minimize the total weighted 
distance between the customers and the closest facilities to them.
%Although unable to develop a method to solve the p-median problem, Hakimi 
%developed the 
%well known ``Hakimi theorem'' that states there exists at least one set of 
%k-medians (one optimal solution) 
%on the nodes of the network (\citet{ReVelle}). 
\citet{ReVelle2} were the first to formulate the $p$-median problem as a 0-1 integer 
program. Since then researchers have developed both exact methods 
(e.g., \citet{Galvao}, \citet{Bilde}) and heuristic methods 
(e.g., \citet{Maranzana}, \citet{Teitz}, \citet{Rosing}) for the $p$-median problem. 
Depending on the difficulty of the instances, optimal or near optimal solutions 
for $p$-median problems are reported for instances up to 3000 customer nodes. 

By removing the constraint that $p$ locations must be selected and by adding a 
fixed cost for each candidate location, we obtain a simple plant location 
problem from the $p$-median problem. Also called the uncapacitated facility 
location problem (UFLP), the problem has been well-studied in the literature 
(\citet{Cornuejols2}). The objective of the problem is to select a subset 
of the candidate facilities and to assign customers to the selected facilities 
so as to minimize the sum of fixed and operating costs. The operating cost in 
the UFLP is the transportation cost calculated assuming individual customer 
routes that originate at a facility, serve one customer and return to the 
same facility. The problem is NP-hard (\citet{Cornuejols2}). A number of 
researchers have focused on polyhedral analysis of the UFLP. Although a 
complete description of the UFLP polytope is not known, \citet{Guignard}, 
\citet{Cornuejols} and \citet {Cho83a,Cho83b} provide partial descriptions. 
Others have focused on developing algorithms for the UFLP, for example, 
branch-and-bound based algorithms by \citet{Erlenkotter} and \citet{Korkel}, 
a Lagrangian relaxation algorithm by \citet{Cornuejols77} and a dual 
simplex algorithm by \citet{Simao}. The best results for the problem are
reported by \citet{Korkel} using an approach that extends the dual-ascent 
heuristic procedure developed by \citet{Erlenkotter} (\citet{ReVelle}). 

%\citet{Morris} and 
%%\citet{ReVelle2} show that linear relaxations of UFLP and p-median problem result 
%integer or near-integer solutions that means, a few branching 
%in a branch and bound algorithm is enough to reach optimal integer solution. 

By adding capacity constraints on the facilities, we obtain the capacitated 
facility location problem (CFLP) from the UFLP.   \citet{Leung} and \citet{Aardal}  
investigate the polyhedral structure of the CFLP and introduce facet defining 
inequalities for the problem, which lead to cutting plane algorithms 
(\citet{Aardal2}) for solving the problem. By using an extended formulation 
(with other valid inequalities) of the CFLP, \citet{Cornuejols} find and compare 
lower bounds for the CFLP by relaxing different constraints within a Lagrangian 
relaxation and decomposition algorithm. Other papers presenting various solution 
algorithms for the CFLP are \citet{Baker}, \citet{VanRoy}, \citet{Pirkul} and 
\citet{Beasley2}. 

In the CFLP, the demand of a customer can be supplied by more than one 
facility. However, in some applications, there may be a single source requirement 
for a customer, which means that the total demand of a customer must be supplied 
by exactly one facility. When the single source requirement is added to the CFLP, 
it is called the capacitated facility location problem with single sourcing 
(CFLPSS). The CFLPSS is more difficult than the CFLP since, while the CFLP is a 
mixed integer model, the CFLPSS is a pure integer model with same number of
variables and constraints. In addition, for a given set of open facilities, 
the CFLP solves a transportation problem to find the assignment of customers to 
the facilities, but the CFLPSS solves a generalized assignment problem, which is 
also an NP-hard problem. In the literature, most of the studies use heuristics, 
especially Lagrangian based heuristics, such as \citet{Klincewicz} and 
\citet{Sridharan}. \citet{Delmaire} design several metaheuristics and 
\citet{Ahuja04} develop a neighborhood search algorithm for the problem.
%Most of the studies use Lagrangean heuristics to solve the problem. 
%Two of these are \citet{Klincewicz} and \citet{Sridharan}.  
\citet{Holmberg} develop an exact branch-and-bound algorithm that uses a  
Lagrangian based heuristic and repeated matching heuristic that solves 
successive matching problems. They can solve problems up to 
200 customers and 30 facilities. 



%[if this is daskinin worku ise koy]
%\citet{Daskin} gives details of two approaches to solve CFLP with lagrangean relaxation algorithm. 
%In the first approach, the capacity constraints are relaxed and what is left is an UFL problem to solve. In the second approach,
%the demand constraints for customers are relaxed and this leaves two subproblems: a continuous
%knapsack problem and a problem closely related with 0-1 knapsack problem. 

%For the sake of LRS problem, we focus on the minimizing cost objective in our literature survey, 
%but in variations of location problems, there are different objectives other than minimizing 
%cost such as minimizing the maximum time to reach a customer or maximizing the number of demands 
%satisfied within a specified service time. Formulation of the problem may change based on the objective.
%Some the basic examples are set covering problems and center problems. For the details of each 
%problem such as their application areas, classifications, formulations and solution methods, 
%see \citet{Daskin} and \citet {Drezner}. 

%[maybe put another place]
%(see the book edited by \citet{Drezner95} for a  survey of applications and methods).
% Considering the variations of 
%location problem, the complexity of the problem changes from simple linear, single-stage,
%single-product, uncapacitated, deterministic models to nonlinear probabilistic models (\citet{Klose})


\subsection{Vehicle Routing}
\label{lit-vrp}
The second component of the LRS problem is how to design the vehicle routes. The 
classical vehicle routing problem (VRP) finds a set of minimum cost vehicle 
routes, each of which starts at the depot, visits some customer locations and 
returns to the depot. There is a large body of literature related to the VRP that 
describes the many problem variations and solution approaches. The VRP has many 
variations depending on the characteristics of the vehicles, the facilities and 
the customers. 
%Vehicle routing problems vary depending on the characterisitics of the vehicles and the facilities. 
%\citet{Bodin} give a taxonomy for VRP. 
For example, the vehicles may be identical or different with respect to capacity 
limits; the problem may be concerned with delivery only, pick up only or both 
delivery and pick up; the vehicles might be restricted to serve each customer in 
a given time window and/or serve according to precedence relations of the customers.
In addition, the problem may involve a single facility or multiple facilities.  
\citet {Christofides} surveys the different types of formulations used to model 
the VRP and the exact and heuristic algorithms designed for the problem prior to 1985. 
In a recent study, \citet{Toth} give a detailed description of the variants of the 
VRP and their formulations and they discuss the solution methods that are applied 
to these variants. 
\citet{LetchfordG06} review the different formulations and derive
relationships between these formulations. Using projection, they compare the LP 
relaxation bounds and show that some valid inequalities in the two index space
are implied by extra flow variables or set partitioning variables.

To formulate the VRP as an integer program, three different approaches are used in 
literature. One approach is to define integer variables for each arc or edge in the 
model representing the number of times the arc or edge is traveled. This type of 
model, which is called a vehicle flow model, can be divided into two groups: 
two index formulations and three index formulations (\citet{Laporte2}). In a 
two index formulation, a flow variable has two indices representing the 
origin and the destination of the arc and in a three index formulation,
each variable has an additional index representing the vehicle traveling the arc.  
Vehicle flow models, especially two index formulations, are very common in the 
VRP literature. \zaupdate{For capacitated VRP (CVRP), \citet{LetchfordG06}
investigate the relationship between two index and three index vehicle flow models
and present some theoretical results. They prove that the 
three index formulation strengthened with certain classes of valid inequalities
gives the same lower bound with a two index formulation with simpler constraints.
Therefore, the LP relaxation bound based on the three index formulation is weaker than
the LP relaxation bound based on the two index formulation.}
In addition to the constraints to form feasible routes, the vehicle 
flow model includes generalized subtour elimination constraints, the number of which 
grows exponentially with the number of nodes in the problem. There is one subtour 
elimination constraint for each subset of customers, which forces a lower bound for 
the number of vehicles exiting the subset based on the minimum number of vehicles 
necessary to serve the customers in the subset. The exponential number of constraints 
is handled by considering a relaxation of the model that includes only some of the 
constraints or by using a branch-and-cut algorithm. \citet*{Miller} introduce a 
polynomial number of constraints that eliminate subtours, but the LP relaxation of the 
model with these new constraints is weaker than the model with generalized subtour 
elimination constraints. \citet{Desrochers3} strengthen the Miller-Tucker-Zemlin 
subtour elimination constraints by lifting them. They define stronger valid 
inequalities that can be adapted to variations of the VRP, such as the capacitated 
VRP (CVRP), the distance constrained VRP (DVRP) and the VRP with time windows (VRPTW).    
  
A second modeling approach, developed by \citet{Garvin}, is called a commodity flow 
model. In this approach, an additional variable for each arc is added to the variables 
in the vehicle flow model. These additional variables represent the amount of demand 
that flows along the edges traveled by the vehicles. \citet{LetchfordG06} 
analyze three different commodity flow models: one commodity, two commodity and multi 
commodity flow models. In two commodity flow model, in addition to the flow on each arc
representing the vehicle load, there is a reverse flow equal to the vehicle residual 
capacity. In multi commodity flow model, demand of each customer represents one commodity
and flow of each commodity is denoted by different variables on each arc. 
\citet{Gouveia95} and \citet{LetchfordG06} state some theoretical results about
the relationship between commodity flow model and two index vehicle flow model.
No exact solution algorithms for 
this approach are reported in the literature. Based on the computational experiments, 
\citet{Gavish} report large gaps between the  LP solutions and the optimal integer 
solutions when the number of customers are greater than 15. 

A third approach, 
proposed by \citet{Balinski}, is a set partitioning formulation, in which there is 
one  variable corresponding to each feasible vehicle route. The model has possibly 
exponentially large number of binary variables representing whether a specific route is 
selected or not. Since the feasibility of each route is considered in the process 
of constructing the set of routes to be included in the model, various routing constraints, 
such as vehicle capacity, route length restrictions and time windows for the customer 
nodes, can be included in the model without changing the formulation. Because of the 
exponentially large number of variables, it is usually solved by a column generation 
algorithm. The LP relaxation 
of the model is generally tight (\citet{Toth}). \citet{LetchfordG06} theoretically
compare the strength of the set partitioning formulation with the two index vehicle 
flow formulation and prove that the set partitioning variables imply some of the
valid inequalities in the two index space.
    
In reviewing the solution algorithms for the VRP, we highlight those for deterministic, 
capacitated VRPs. In a capacitated VRP (CVRP), the total demand of a route must be less than 
the capacity of the vehicle. Both exact methods and heuristic methods have been developed 
to solve the CVRP. \citet{Laporte2} survey the exact algorithms for solving the VRP. 
They give a complete analysis of branch-and-bound algorithms, which are the most common 
way to solve the CVRP. They also describe dynamic programming and integer programming 
approaches that use different formulations of the problem. 
In addition to branch-and-bound procedures, 
\citet{Laporte3} describes set partitioning formulations, column 
generation procedures and heuristic approaches (e.g. Clark and Wright, tabu search 
and sweep algorithm) for the problem. 

Branch-and-cut algorithms are very successful in finding optimal solutions for large 
instances of the symmetric TSP, which is a problem closely related to the CVRP. This 
success has motivated the development of branch-and-cut algorithms for vehicle flow 
models of the CVRP. \cite{Toth} describe how to apply a branch-and-cut algorithm to the 
VRP and the CVRP. Branch-and-cut algorithms incorporate the results of polyhedral studies 
of the CVRP polytope, valid inequalities (e.g. \citet{Laporte4} and \citet{Laporte5}) and 
facets of the polytope. The first papers published about the facets of the CVRP polytope 
were by \citet*{Araque1} and \citet{Araque2}, who  assumed unit demand for customers, and 
by \citet{Cornuejols93}, who extended the research to general demand. One of the main 
difficulties with branch-and-cut algorithms is that the separation procedure for the cuts 
might be NP-hard. To overcome this difficulty, researchers seek polynomial time
algorithms for exact separation (e.g., \citet{McCormick} and \citet{Blasum}) and
heuristic separation algorithms (e.g., \citet{Augerat2}, \citet{Ralphs}). 
%Heuristic separation approaches (e.g., \citet{Augerat2}, \citet{Ralphs}) 
%are designed to overcome this difficulty. For capacity constraints in the CVRP, \citet{Toth} 
%note that \citet{McCormick} and \citet{Blasum} define polynomial time algorithms for the 
%exact separation.   
 
%For set covering or set partitioning-based models of the CVRP, a number of different 
%approaches are used. 
Set covering-based formulations are used successfully by \citet{Cullen} 
to solve the VRP heuristically and by \citet{Agarwal}, \citet{Bixby}, \citet{Hadjiconstantinou} 
and \citet{Desrochers2} to solve the VRP exactly. The first step for solving a set 
partitioning-based formulation is to solve the LP relaxation of the model using column 
generation. The objective of the column generation subproblem (pricing problem) is 
to construct feasible vehicle routes which include each customer at most once 
(elementary routes) and satisfy some other constraints required for the feasibility 
of a route such as vehicle capacity. Since the complexity of the problem of finding 
elementary routes is NP-hard, it is common to relax this property and solve the set
partitioning formulation with non elementary routes.

%%%%% ************************* Bi tuhaflik var bu paragrafta *****************************
%%%%% ???????????????????????????????????????????????????????????????????????? 
\zaupdate{Read Toth sections and fix the following paragraph. There seems
to be something wrong}

\citet{Agarwal} use a branch-and-bound algorithm, \citet{Desrochers2} use a 
branch-and-bound algorithm and dynamic programming, \citet{Bixby} uses a cutting plane 
algorithm and \citet{Hadjiconstantinou} solve the dual of column generation problem using a 
dynamic programming heuristic to get the LP relaxation of the problem. For many of the 
problems, the LP relaxation solution is very close to integer solution. However, even if 
the LP relaxation is very strong, some branching may be required to obtain an integer 
solution. Both cutting plane algorithms  (see \citet{Toth} for details) and 
branch-and-price algorithms (see \citet{Desrochers2}) have been used. 
%[look at problems that use network problems to solve column generation] 

\zaupdate{\citet{Fukasawa06} combine cutting plane and column generation approach
by extending the set partitioning formulation for CVRP with
the two index vehicle flow variables and constraints based on these variables. They
design a branch, cut and price algorithm which seems to be very successful compared to
the reported results in the literature.} 

A generalization of the VRP, which is also a special case of our LRS problem, is the Multi-Depot 
Vehicle Routing Problem (MDVRP). The MDVRP constructs a set of vehicle routes with 
minimum cost in which each customer is visited by exactly one vehicle and 
each vehicle route starts and ends at the same depot. Exact solution algorithms are very 
rare for MDVRP; \citet{Laporte6} and \citet{Laporte7} are the only two papers that report 
exact solution algorithms for the MDVRP. Both solution methods are based on branch-and-bound 
algorithms. Many papers develop heuristic algorithms based on construction 
and improvement techniques (e.g., \citet{Tillman}, \citet{Tillman2}, \citet{Gillet} and 
\citet{Chao}) and  tabu search (e.g., \citet{Cordeau} and \citet{Renaud}).

%look at literature of renaud and literature of MDVRP inter depot trips by clevier

\subsection{Integrated Facility Location and Routing}

In most facility location problems, the assumption is that customers are served on individual routes 
or that customers travel to the facilities. For situations in which customers receive service from 
routes making multiple stops, the assumption will not accurately capture the transportation cost.
In order to accurately represent the cost of serving customers, we need to know the order of the 
customers on the route. Hence, the location and routing problems must be solved simultaneously. 
As early as 1968, \citet{Webb} realized that using straight line trips to each customer versus
using multiple stop tours resulted in different solutions. \citet{Salhi} also investigated the 
effect of ignoring the routing of vehicles while deciding the locations of facilities. They show 
that solving the location and routing problems independently does not always result in  the best 
solution in terms of cost. Thus, although solving the combined problem may be more difficult than 
solving the individual ones, it may result in a more cost efficient solution. 

%As the distance of each customer to each facility is a major
%concern in the formulation, the routing of the vehicle while serving
%customers is ignored. In transportation networks, because of the
%interdependency between vehicle routes and the location of the facilities,
%optimization of these decisions simultanesouly might result less cost and
%more efficient solution than the solution of location and routing problems
%indiviually.  Indeeed, 
%[Webb MHJ. Cost functions in the location of depot for multiple delivery journeys. Operational Research Quarterly
%    1968;19:311 ­ 20.

\citet{Laporte} reviews the work on deterministic location routing problems; 
he describes the different formulations of the problem, the solution methods 
used and the computational results published up to 1988. A more recent paper 
by \citet{Min} classifies both deterministic and stochastic problems in the LRP 
literature with respect to problem characteristics and solution methods. They 
group the LRP papers according to characteristics such as single or multistage, 
deterministic or stochastic, single or multiple facilities and vehicles, and  
uncapacitated or capacitated vehicles and facilities. Because the LRP combines two NP-hard
problems, researchers have developed a wide variety of both heuristic and exact algorithms for 
the LRP. Most of the heuristic algorithms tend to divide the problem into subproblems. 
Heuristic algorithms handle these subproblems sequentially or iteratively. Based on \citet{Min2}, 
for problems with large facility fixed costs and capacitated vehicles, sequential approaches perform 
computationally better. Some examples of sequential heuristic algorithms are 
developed by \citet{Perl}, \citet{Hansen} and \citet{Barreto07}.
\citet{Perl} develop a heuristic for a three index flow formulation of an LRP 
with capacitated distribution centers and 
vehicles. They decompose the problem into three subproblems: a multi-depot vehicle 
dispatch problem, a warehouse location-allocation problem and a multi-depot routing 
allocation problem. They develop an algorithm 
that solves the problems sequentially but considers the dependency between the problems.  
They solve the  multi-depot vehicle dispatch problem and the multi-depot routing 
allocation problem heuristically and the location-allocation problem exactly. 
\citet{Hansen} introduce flow variables and flow constraints to the model in 
\citet{Perl} to get an MILP and solve the new model with a heuristic that is a 
modification of the heuristic developed by \citet{Perl}. \citet{Barreto07} develop 
a sequential heuristic algorithm for problems with capacitated facilities and
capacitated vehicles. In their heuristic, the first step is to solve routing 
subproblem. They group customers using clustering techniques and design TSP solutions. 
Then, by assuming each group of customers as a single customer, single source 
capacitated location routing problem is solved. \citet{Srivastava} presents 3 iterative 
heuristics called SAV1, SAV2 and CLUST and compares the results with a classical 
sequential heuristic in which facility locations are determined first and 
vehicle routes are designed based on open facilities. SAV1 heuristic starts with all 
depots being open and using a multi depot vehicle routing heuristic calculates 
approximate routing costs for the facilities. At each step a facility is closed 
until the problem becomes infeasible. SAV2 heuristic works similar to SAV1 except 
initially all depots are closed and at each iteration one depot is opened. CLUST
heuristic generates groups of customers using a minimum spanning tree and a density search 
clustering technique. Then a depot is selected for each cluster considering the 
distance between cluster centroid and the depot.

\zaupdate{And finding LB in Albareda Sambola }


Exact algorithms for the LRP have been developed mostly for two index vehicle flow formulations. 
\citet{Laporte81} solve a single depot model with a constraint relaxation method
and a branch-and-bound algorithm. They report solving problems with 50 customer locations.
\citet{Laporte83} solve a multi-depot model using a constraint
relaxation method and Gomory cutting planes to satisfy integrality. They can solve problems
with at most 40 customer sites. \citet{Laporte86}
apply a branch-and-cut algorithm to a multi-depot LRP model with vehicle capacity constraints. They
use subtour elimination constraints and chain barring constraints that guarantee that each route
starts and ends at same facility. The authors report computational results for problems with 20
customer locations and 8 depots. \citet{Belenguer06} provided a two-index formulation of
the LRP with capacitated facilities and capacitated vehicles and
presented a set of valid inequalities for the problem. They developed
two branch-and-cut algorithms based on different formulations of the
problem. They reported that they could solve instances with up to
32 customers to optimality in less than 70 CPU seconds and could provide
good lower bounds for the rest of the instances, which had up to 134
customers.
%Examples in the 
%literature include a single depot model by \citet{Laporte81}, a multi-depot model by \citet{Laporte83} 
%and a multi-depot model with vehicle capacity constraints by \citet{Laporte86}. 
\citet{Berger} develops a set partitioning-based model for an LRP with route-length constraints and 
uncapacitated vehicles and facilities. \citet{Berger05} extend the work to develop an exact 
branch-and-price algorithm in which they solve the pricing problem as an elementary shortest 
path problem with one resource constraint. They report computational results for problems with 
100 customers and various distance constraints. 



\subsection{Vehicle Scheduling}

The third component of the LRS problem is how to assign routes to vehicles, which is known 
as the vehicle scheduling or loading problem. The (multi-depot) vehicle scheduling problem 
((MD)VSP) seeks a least-cost assignment of routes (or trips), each with a starting and 
ending time, to vehicles such that each route is assigned to one vehicle, each vehicle 
starts a route after the starting time and finishes before the ending time and each vehicle 
starts from a depot and returns to same depot. In some variations, there are constraints 
such as depot capacity and upper and lower bounds on the number of vehicles at each depot. 
The MDVSP is a special case of the LRS problem; given a set of selected facilities and a 
set of vehicle routes connected to these facilities such that each customer is served by 
one vehicle, the LRS problem reduces to a MDVSP.  \citet{Carraresi} describe how to convert 
a single depot VSP to an assignment problem or a minimum cost network flow model 
and therefore they conclude that the single depot VSP is easy to solve. However, \citet{Bertossi} 
show that if there are more than 2 depots, the MDVSP is NP-hard. 
 
Both heuristic and exact solution methods are proposed for the MDVSP. \citet{Bodin} 
propose two heuristics based on decomposing the problem into single depot VSPs.  
\citet{Bertossi} model the MDVSP by defining an independent problem for each depot and 
additional constraints to cover each trip exactly once. They relax the additional constraints 
using a Lagrangian relaxation algorithm and then solve independent VSPs. 
\citet{DellAmico} propose a heuristic method that guarantees a solution with the minimum 
number of vehicles. They solve shortest path problems on a graph that includes one 
node for each route and arcs between two nodes that can be assigned to the same 
vehicle. A path on the graph represents a vehicle duty.  
 
Several exact algorithms based on different formulations of the problem are 
proposed. \citet{Carpaneto} use a branch-and-bound algorithm based on additive lower 
bounds. \citet{Forbes} develop an algorithm using a dual simplex method to solve their 
three-index integer programming model, which is a multi-commodity network flow problem. 
\citet{Ribeiro} and \citet{Bianco} use set partitioning-based formulations for which 
they generate columns and apply a branch-and-bound algorithm to get an integer solution. 
\citet{Fischetti} design a branch-and-cut algorithm by defining valid inequalities of the MDVSP.  

%[an option]
%\citet{Bodin} propose two different approaches based on decomposing the problem into simpler problems. The 
%first approach, called cluster first - schedule second, partitions the trips  into sets for each depot and 
%solves a single depot VSP. The second approach, called schedule first - cluster second, solves single depot 
%VSP and then assign the vehicle duties (a duty is the set of trips assigned to a vehicle) to the depots by 
%considering capacity constraints and the minimum cost objective. \citet{Bertossi} model the MDVSP by 
%defining an independent problem for each depot and additional constraints to cover each trip exactly once. 
%They relax the additional constraints using a Lagrangean relaxation algorithm and then solve independent 
%VSPs. \citet{DellAmico} propose a heuristic method that guarantees a solution with minimum number of vehicles. 
%They construct a graph including one node for each route and arcs between two nodes that can be assigned to 
%the same vehicle. At each step of the algorithm, they choose a set of arcs as forbidden arcs and find a 
%shortest path avoiding these arcs to determine a duty of a vehicle.



%[OLD LONG VERSION]
%\citet{Carraresi} describe how to convert a single depot VSP to a network model. They construct a bipartite 
%network with disjoint node sets of S and T. Both sets S and T include a node representing to a trip or a 
%task that is assigned to a vehicle.  Arcs in the bipartite graph  are created based on scheduling restrictions 
%for the trips. They prove that there is a one-to-one correspondence between vehicle schedules and matchings 
%in the constructed bipartite graph. In addition, by adding some arcs and nodes, they construct another network 
%in which they prove that each integral feasible flow or assignment corresponds to a feasible vehicle schedule. 
%Furthermore, they can include costs for arcs to handle different objectives. 
%Finally, they conclude that since a vehicle scheduling problem can be converted to an assignment or min 
%cost flow problem, it is easy to solve a VSP. However, if there are more than 2 depots, the MDVSP is proved to 
%be NP-Hard (\citet{Bertossi}). Therefore, heuristics were proposed in the earliest MDVSP papers.

%\citet{Bodin} propose two different approaches based on decomposing the problem into 
%simpler problems. The first approach, called cluster first - schedule second, partitions the trips  into sets 
%for each depot and solves a single depot VSP. If the solution is not feasible with respect to the capacity of 
%the facility, the clusters are updated and the VSP is resolved. The second approach is schedule first - cluster 
%second. They solve a single depot VSP and then assign the vehicle duties (a duty is the set of trips assigned 
%to a vehicle) to the depots by considering capacity constraints and the minimum cost objective. \citet{Bertossi} 
%model the MDVSP by defining an independent problem for each depot and additional constraints to cover each trip 
%exactly once. They relax the additional constraints using a Lagrangean relaxation algorithm and then solve 
%independent single depot vehicle scheduling problems. Later, \citet{DellAmico} propose a heuristic method that 
%guarantees that the minimum number of vehicles is used. They construct a graph including one node for each trip 
%and arcs between nodes if the trips can be assigned to the same vehicle. At each step of the algorithm, they 
%choose a set of arcs as forbidden arcs and find a shortest path avoiding these arcs to determine a duty of a 
%vehicle.  
       
%Several exact algorithms based on different formulations are proposed in literature. \citet{Carpaneto} formulate 
%the problem as an integer linear programming model and solve the problem using a branch and bound algorithm 
%based on additive lower bounds. Their formulation includes assignment constraints and flow conservation 
%constraints to form the paths. To get lower bounds, they combine the bounds obtained from the assignment and the 
%shortest path relaxations of their model. \citet{Forbes} formulate a three-index integer programming model that 
%is a multi-commodity network flow problem. They solve the problem using the dual simplex method. \citet{Ribeiro} 
%formulate a set-partitioning based model and use column generation to find lower bounds. They develop a branch 
%and bound algorithm to get optimal solutions. \citet{Bianco} also use a set-partitioning based formulation to 
%model the MDVSP. They generate a good subset of duties by heuristically solving the dual of the LP relaxation 
%of the model. Finally, they use branch-and-bound to find the optimal solution. \citet{Fischetti} define some 
%valid ineqaulities of the MDVSP and design a branch-and-cut algorithm. They get near optimal solutions within 
%their time limits.
 
\subsection{Vehicle Routing and Scheduling}
\label{lit-vrps}
In the VRP, there is one-to-one relationship between vehicles and routes, 
which means that one route is assigned to each vehicle. In a combined vehicle 
routing and scheduling problem (VRSP), scheduling constraints
are added to the VRP constraints to change this one-to-one relationship between 
vehicles and routes. Therefore, the VRSP seek to determine the set of vehicle 
routes as well as the assignment of these routes to the available vehicles. 
In the literature, this version of the VRSP is sometimes called the VRP with 
multiple trips (VRPMT). Only a few papers in the literature address the 
VRPMT and they only describe heuristic approaches.  

\citet{Taillard} use a three-step algorithm to find good feasible solutions 
for the VRSP. In the first step, they generate some number of VRP solutions 
using a tabu search algorithm. In the second step, they select a subset of 
these routes using an enumerative algorithm and, in the third step, they use 
a bin packing heuristic to combine the routes in such a way that they will 
not violate the time constraint and assign them to vehicles. The authors 
generate test problems with 50 to 199 customers. Feasible solutions are 
provided for 82 instances out of 104 instances. They report that the cost 
of the feasible solutions are on the average  within 1.2\%  of the VRP 
solution obtained by \citet{Rochat}.     

\citet{Brandao} use a tabu search algorithm and the GENI algorithm 
(developed by \citet{Gendreau} to solve TSP) to solve the VRPMT with time 
window constraints. They design a three-phase algorithm in which the 
solution is improved in each phase. The input for the first phase is an 
initial set of routes constructed using VRP constraints.
%and input of the a step is created by the previous step. 
The first phase checks each route in the initial set for the time window 
infeasibility and, if there is an infeasible route, the route is divided 
into smaller routes that satisfy the time windows. The second phase 
constructs a solution with smaller cost by searching the set of routes. 
Similar to the second phase, the third phase also seeks a solution 
with lower cost but the solution space is restricted to the feasible 
space during the search. They evaluate their algorithm by comparing their 
solutions with manual solutions in a case study. They report on average 
20\% better solutions. In addition, \citet{Brandao2} solve the instances 
created by \citet{Taillard}. They can find feasible solutions for 89 
instances out of 104 instances.     

\citet{Petch} develop a multi-stage heuristic algorithm in which they 
create vehicle routes and then combine the routes using bin packing 
heuristics. \citet{Olivera} solve the VRPMT using an adaptive memory 
algorithm that is an enhancement of tabu search. The basic idea is to 
construct new solutions using the components of the previously generated 
solutions. They initialize the memory by generating some feasible 
solutions to the VRP, which they use to construct feasible assignments 
of routes to vehicles. Then, they use tabu search to improve the 
constructed solution. This heuristic differs from the others in that 
it constructs feasible vehicle assignments in each iteration, 
while the others only construct the assignment at the final stage. 
They can find feasible solutions for 95  instances out of 
104 \citet{Taillard}'s instances. The feasible solutions are on the 
average within 1.6\% of the best known VRP solutions. In infeasible 
solutions, in which overtime is used, the solutions provided by 
\citet{Olivera} are better than the solutions provided by other authors.      



%Besides scheduling of vehicles, if the problem also seeks for the vehicle routes, then it is called
%combined vehicle routing and scheduling problem (VRSP) and the problem is NP-hard (\citet{Lenstra}). 
%The most common VRSP is vehicle routing problem with time windows (VRPTW). In the problem, in addition 
%to the VRP constraints, each customer has a time interval and must be visited in its time interval. 

%\citet{Toth} give detailed desription of the formulations, solution methods used and computational results obtained 
%for VRPTW. They present both multicommodity flow model with three index flow variables and set partitioning
%based model. They survey heuristics algorithms proposed and used to solve VRPTW such as heuristics based on route 
%contruction, route improvement, combination of these two, metaheuristics such as tabu search, simulated annealing,
%and genetic search. Decomposition approaches such as Lagrangean relaxation and Danzig Wolfe decomposition
%algorithms are very common to solve VRPTW. Danzig Wolfe decomposition model is equivalent to set partitioning based 
%model. 

%[effectiveness of set covering formulations on VRPTW]

%\citet{Desrochers2} solve VRPTW by formulating the problem as a set partitioning model that is equivalent to a set 
%covering prolem when the triangular inequality is satisfied for the distance matrix. They use column generation
%algorithm to solve the LP relaxation of their model. A dynamic programming algorithm is used to solve the subproblem
%that generates columns. The exact subproblem to create columns is an elementary shortest path problem with vehicle 
%capacity and time window constraints (ESPPTWCC). Ignoring the elementary property of a path, they seek shortest paths that
%does not include 2-cycles. They use branch and bound algorithm to get optimal or a good upper bound for the 
%optimal solution. First, they branch on total number of vehicles used and total distance travelled (if distance matrix is 
%integer), second they branch on the arcs in the network. A score is assigned to each arc based on the flow on the arc. 
%While choosing arcs to branch on based on their scores, they first consider columns that have cycles, then they 
%pick arcs from fractional variables. They can solve problems up to 100 customers. 

%\citet{Kohl} solves VRPTW using a Lagrangean relaxation algorithm in which they relax the constraints
%that say that each customer node must be visited. When they relax these constraints, the remaining 
%problem is elementary shortest path problem with time window and capacity constraints (ESPPTWCC). For any lagrangian
%multiplier vector, their subproblem problem decomposes into equivalent ESPPTWCC for each vehicle. Their master problem 
%is to find the optimal Lagrangean multipliers that maximizes the lower bound of the VRPTW. Similar to 
%the \citet{Desrochers2}, they solve relaxation of ESPPTWCC, shortest path problem with constraints. 
%They use several ways to get better paths such as adding new nodes to the network in order to eliminate negative cycles
%and solve the shortest path with no 2-cycles. They solve their master problem by using both subgradient optimization
%and bundle method. When they find the optimal lagangean multipliers, they do branch and bound algorithm to get the
%optimal integer solution. They branch on time windows for the customers by dividing the time window into two and 
%using these new smaller windows in new nodes. They can solve the problems solved by \citet{Desrochers2} and also some other
%unsolved problems.    
   
%In VRPTW, there is one-to-one relation between vehicles and trips which means that one vehicle is assigned to one
%route, but in practical applications, the fleet size might be small or while the capacity of trucks are small, 
%the working hour limit is high enough to schedule more than one trip to one vehicle. This extension of VRP is 
%actually a VSP and in literature it is called VRP with multiple trips (VRPMT). The study of VRPMT is very limited
%in literature and furthermore, there is not any paper in literature that we are aware of and that solve the exact VRPMT.  
   

%\citet{Bodin} define and classify the vehicle scheduling problems. According to their definition, if the vehicles
%must visit customer nodes at some specified times, then the problem is called as vehicle scheduling problem. 
%If the problem also must decide the order of customer visits, then the problem is called the combined vehicle routing
%and scheduling problem. The basic combined vehicle routing and scheduling problem seeks to minimize the total
%number of vehicles and the total cost of routing and serving to customers subject to
%scheduling constraints and capacity restrictions.
%In practical problems, there might be a time window for serving each customer and restrictions for the driver's
%working hours. Because integer programming formulations are difficult to solve
%for large instances, heuristic methods are prefered for practical size problems. 

%look-Lenstra and Kan ,{C}omplexity of {V}ehicle {R}outing and {S}cheduling {P}roblems,{1981},{Networks},{11}, {221-228}



\subsection{Integrated Location, Routing and Scheduling}

The decisions of location and routing are interdependent for the cases in which customers are served
by routes with multiple stops. In such cases, solving the combined LRP may result in a more cost 
efficient solution. We can go one step further and say that most LRP papers assume that there is 
a one-to-one relationship between routes and vehicles.  
However, it may be that the vehicle capacity restricts the route but the total available 
time does not. In this case, a vehicle can serve more than one route, and thus, the decision 
of assigning of routes to vehicles should be optimized together with the decisions of facility 
location and vehicle routing. Integrating these three decisions gives rise to a class of problems 
called {\em integrated location routing and scheduling} (LRS) problems.  

As far as we know, the papers by \citet{Lin} and \citet{Lin2} are the only LRS papers to appear in the 
literature. \citet{Lin} describe an LRS problem, which arises in the delivery of telephone bills, in 
which they seek to locate facilities and to route and assign vehicles so as to minimize cost. 
To solve the problem, they develop a heuristic that divides the problem into three phases, facility 
location, routing and loading, and that iterates through the phases. The location phase determines the set 
of facilities via enumeration. For a given set of opened facilities, the routing phase constructs vehicle 
routes using the Clarke and Wright algorithm (\citet{Clarke}) followed by an insertion heuristic and 
more extensive metaheuristics. The loading phase optimally assigns routes to vehicles considering the 
time limits. The authors report the development of a branch-and-bound algorithm to solve the exact problem 
to evaluate their heuristics but they do not provide a formulation. They test their algorithm on five 
instances with 10-12 customers and 4 potential depots, but the branch-and-bound 
algorithm can solve only three of the five instances to optimality within the specified time 
limit. For these instances, they report that their heuristic algorithm can find optimal or good solutions 
quickly. In addition, they solve a real life instance with 27 customers and 4 depots using their
heuristic algorithm and compare the results with the management's manual solution. 


\citet{Lin2} extend the study by \citet{Lin} to solve a multi-objective LRS problem that considers
workload balance in terms of total vehicle time and total vehicle load as well as total cost. They
use the same three-phase heuristic algorithm designed by \citet{Lin}. Since they consider multi-objective
problems, for each feasible solution in the algorithm, they keep three performance measures, the total cost,
the difference between the maximum and the minimum vehicle time and the difference between maximum and minimum
vehicle load. The feasible solutions are evaluated based on these three performance measures and only non-dominated 
solutions are kept. In addition, they design a tabu search (TS) algorithm  for the routing phase of the algorithm and 
compare its performance with the simulated annealing (SA) algorithm applied in \citet{Lin}. While the TS algorithm
and the SA algorithm have similar performance in some instances, the SA algorithm is better than the TS algorithm
when the number of demand nodes in a route increases. Furthermore, 
they compare the performance of the algorithm in two approaches: assigning vehicles to routes after 
the routing phase (sequential approach) and assigning vehicles while constructing routes
(multiple approach). They conclude that the performance of the approaches depends on the difficulty
of the instance. The multiple approach is better than the sequential approach when the solution includes
routes with more demand nodes. 
They test their algorithm with two sets of real data, one with 27 customers and one with 89 customers, 
and a set of simulated data including instances with 100 and 200 customers. 


%The earliest paper similar to LRS is in 1981 by \citeauthor{Nambiar}. He formulates a 
%problem in the natural rubber industry that includes capacity restrictions for the facilities 
%and the vehicles and time constraints on the tours. They formulate a mixed integer program that 
%seeks the optimal location of facilities, assignment of stations to these 
%facilities and the number of vehicles as well as their routes. They propose two 
%different heuristics based on the decomposition of the problem into a facility location part 
%and a vehicle routing part. In both heuristics, they group the customers according to the 
%proximity to the facilities, vehicle capacity and time restriction. Then they 
%use a TSP to find routes for the customers in each group. In one heuristic, they 
%select the facilities to be open according to minimum cost. In the other heuristic,
%they solve a simple plant location model. Finally, they use a savings heuristic to 
%improve the routes in both heuristics.    

\subsection{Summary of Literature}

The LRS problem integrates three NP-hard problems: facility location problem, 
vehicle routing problem and vehicle scheduling problem. Solving these three
problems interdependently and simultaneously results in more realistic 
and cost efficient solutions. However, the combined LRS problem is more difficult 
than its components. In the literature, there are only two papers investigating the
LRS problem. These papers only develop heuristic algorithms, they do not provide
a formulation of the problem, an exact algorithm to find the optimal solution and  
a lower bound for the optimal solution.

The LRS problem is a generalization of other NP-hard problems such as LRPs, MDVRPs, VRPMTs. 
There are only few papers solving those problems exactly. For the multi-depot
LRP, the optimal solutions are reported for instances with uncapacitated vehicles and up to 
40 customers and instances with capacitated vehicles and up to 20 customers. For the MDVRP, 
computational results for exact algorithms are provided for instances with up to 50
customers and 8 depots (with symmetric distance matrix) and  for instances with up to 80
customers and 3 depots (with asymmetric distance matrix). For the VRPMT, no exact algorithm
is provided in the literature. Currently, we can solve instances up to 40 customers and 5 facilities.
%Our goal is to solve instances up to 50 customers and 10 facilities.

%One way to solve the LRS problem is to create smaller problems by fixing some of the decisions. 
%We can constuct the LRS problem in three ways in terms of subproblems each of which is also NP-hard. 
%First, for a given set of facilities and vehicle routes, the problem solves a MDVSP. Second, for 
%a given set of facilities and with no time constraint the problem solves a MDVRP. Third, for a given 
%set of vehicle assignments, the LRS problem needs to solve a CFLPSS.  In this proposal,
%we solve the LRS problem using a set partitioning-based model, which is a CFLPSS. We evaluate 
%all possible set of vehicle assignments and solve the CFLPSS model to get an optimal solution
%to the LRS problem.   
       




